% GENERATED FILE. Edit canonical lessons, then run npm run generate:books.
% canonical-source-hash: fnv1a64:5ce87824a9d86515
% canonical-lessons: MR-C225-patavne, MR-C225-tika-karne, MR-C225-siddha-karne, MR-C225-avlambun-asne, MR-C225-talne

\chapter{To convince, To criticise, To prove, To depend, To avoid}
\label{ch:mr-to-convince-to-criticise-to-prove-to-depend-to-avoid}
\hlchaptermodality{225}

\begin{chapteropening}
\textbf{By the end of this chapter:} \emph{I can say to convince, to criticise, to prove, to depend and to avoid.}

Complete the last lesson of chapter 225: I can say to convince, to criticise, to prove, to depend and to avoid.
\end{chapteropening}

\section[\texorpdfstring{\mr{पटवणे} (paṭavṇe)}{पटवणे (paṭavṇe)}]{\mr{पटवणे} (paṭavṇe) — to convince}

\label{lesson:MR-C225-patavne}



\begin{quote}
Before the new one: say the Marathi for to rest, then the Marathi for to get bored.
\end{quote}

\subsection*{You'll want to know: \mr{पटवणे}}
\textbf{\mr{पटवणे}} — \emph{paṭavṇe} — \textquotedblleft{}to convince\textquotedblright{}.

\begin{grammarlens}[title={What you've built}]
One more word for this chapter.
\end{grammarlens}

\subsection*{Your turn}
\begin{itemize}
\raggedright
  \item \emph{Say it:} \emph{paṭavṇe}
  \item \emph{Say it:} \emph{paṭavṇe}, once more
  \item \emph{Say it:} say \emph{paṭavṇe}
  \item \emph{Recall:} say the Marathi for to rest, then the Marathi for to get bored, then say \emph{paṭavṇe} again
\end{itemize}

\begin{tcolorbox}[breakable,colback=teal!4,colframe=teal!35!black,title={Before you move on}]
What does \mr{पटवणे} mean? (\textquotedblleft{}To convince\textquotedblright{}.) Say it once more.
\end{tcolorbox}
\section[\texorpdfstring{\mr{टीका} \mr{करणे} (ṭīkā karṇe)}{टीका करणे (ṭīkā karṇe)}]{\mr{टीका} \mr{करणे} (ṭīkā karṇe) — to criticise}

\label{lesson:MR-C225-tika-karne}



\begin{quote}
Before the new one: say the Marathi for to get bored, then the Marathi for to convince.
\end{quote}

\subsection*{You'll want to know: \mr{टीका} \mr{करणे}}
\textbf{\mr{टीका} \mr{करणे}} — \emph{ṭīkā karṇe} — \textquotedblleft{}to criticise\textquotedblright{}.

\begin{grammarlens}[title={What you've built}]
One more word for this chapter.
\end{grammarlens}

\subsection*{Your turn}
\begin{itemize}
\raggedright
  \item \emph{Say it:} \emph{ṭīkā karṇe}
  \item \emph{Say it:} \emph{ṭīkā karṇe}, once more
  \item \emph{Say it:} say \emph{ṭīkā karṇe}
  \item \emph{Recall:} say the Marathi for to get bored, then the Marathi for to convince, then say \emph{ṭīkā karṇe} again
\end{itemize}

\begin{tcolorbox}[breakable,colback=teal!4,colframe=teal!35!black,title={Before you move on}]
What does \mr{टीका} \mr{करणे} mean? (\textquotedblleft{}To criticise\textquotedblright{}.) Say it once more.
\end{tcolorbox}
\section[\texorpdfstring{\mr{सिद्ध} \mr{करणे} (siddha karṇe)}{सिद्ध करणे (siddha karṇe)}]{\mr{सिद्ध} \mr{करणे} (siddha karṇe) — to prove}

\label{lesson:MR-C225-siddha-karne}



\begin{quote}
Before the new one: say the Marathi for to convince, then the Marathi for to criticise.
\end{quote}

\subsection*{You'll want to know: \mr{सिद्ध} \mr{करणे}}
\textbf{\mr{सिद्ध} \mr{करणे}} — \emph{siddha karṇe} — \textquotedblleft{}to prove\textquotedblright{}.

\begin{grammarlens}[title={What you've built}]
One more word for this chapter.
\end{grammarlens}

\subsection*{Your turn}
\begin{itemize}
\raggedright
  \item \emph{Say it:} \emph{siddha karṇe}
  \item \emph{Say it:} \emph{siddha karṇe}, once more
  \item \emph{Say it:} say \emph{siddha karṇe}
  \item \emph{Recall:} say the Marathi for to convince, then the Marathi for to criticise, then say \emph{siddha karṇe} again
\end{itemize}

\begin{tcolorbox}[breakable,colback=teal!4,colframe=teal!35!black,title={Before you move on}]
What does \mr{सिद्ध} \mr{करणे} mean? (\textquotedblleft{}To prove\textquotedblright{}.) Say it once more.
\end{tcolorbox}
\section[\texorpdfstring{\mr{अवलंबून} \mr{असणे} (avlambūn asṇe)}{अवलंबून असणे (avlambūn asṇe)}]{\mr{अवलंबून} \mr{असणे} (avlambūn asṇe) — to depend}

\label{lesson:MR-C225-avlambun-asne}



\begin{quote}
Before the new one: say the Marathi for to criticise, then the Marathi for to prove.
\end{quote}

\subsection*{You'll want to know: \mr{अवलंबून} \mr{असणे}}
\textbf{\mr{अवलंबून} \mr{असणे}} — \emph{avlambūn asṇe} — \textquotedblleft{}to depend\textquotedblright{}.

\begin{grammarlens}[title={What you've built}]
One more word for this chapter.
\end{grammarlens}

\subsection*{Your turn}
\begin{itemize}
\raggedright
  \item \emph{Say it:} \emph{avlambūn asṇe}
  \item \emph{Say it:} \emph{avlambūn asṇe}, once more
  \item \emph{Say it:} say \emph{avlambūn asṇe}
  \item \emph{Recall:} say the Marathi for to criticise, then the Marathi for to prove, then say \emph{avlambūn asṇe} again
\end{itemize}

\begin{tcolorbox}[breakable,colback=teal!4,colframe=teal!35!black,title={Before you move on}]
What does \mr{अवलंबून} \mr{असणे} mean? (\textquotedblleft{}To depend\textquotedblright{}.) Say it once more.
\end{tcolorbox}
\section[\texorpdfstring{\mr{टाळणे} (ṭāḷṇe)}{टाळणे (ṭāḷṇe)}]{\mr{टाळणे} (ṭāḷṇe) — to avoid}

\label{lesson:MR-C225-talne}



\begin{quote}
Before the new one: say the Marathi for to prove, then the Marathi for to depend.
\end{quote}

\subsection*{You'll want to know: \mr{टाळणे}}
\textbf{\mr{टाळणे}} — \emph{ṭāḷṇe} — \textquotedblleft{}to avoid\textquotedblright{}.

\begin{grammarlens}[title={What you've built}]
One more word for this chapter.
\end{grammarlens}

\subsection*{Your turn}
\begin{itemize}
\raggedright
  \item \emph{Say it:} \emph{ṭāḷṇe}
  \item \emph{Say it:} \emph{ṭāḷṇe}, once more
  \item \emph{Say it:} say \emph{ṭāḷṇe}
  \item \emph{Recall:} say the Marathi for to prove, then the Marathi for to depend, then say \emph{ṭāḷṇe} again
\end{itemize}

\begin{tcolorbox}[breakable,colback=teal!4,colframe=teal!35!black,title={Before you move on}]
What does \mr{टाळणे} mean? (\textquotedblleft{}To avoid\textquotedblright{}.) Say it once more.
\end{tcolorbox}
